% !TEX root =  Lecture17.tex


%%%%%%%%%%%%%%%%%%%%%%%%%%%%%%%%%%%%
% Recap/Agenda
%%%%%%%%%%%%%%%%%%%%%%%%%%%%%%%%%%%%
\begin{frame}
\frametitle{Module 4: Bayesian Statistics}
    \begin{itemize}
        \item \hl{Previously:} logistic regression inference (Lecture 16), the last topic of Module 3. Everything so far has been \emph{frequentist}.
        \item \hl{This time:} a new module, and a different way to reason about uncertainty.
          \begin{itemize}
              \item probability review: conditional probability and \hl{Bayes' rule};
              \item putting a \hl{prior distribution} on an unknown parameter, and reading off the posterior;
              \item \hl{sequential updating}, and why the order of the data doesn't matter;
              \item how the Bayesian and frequentist answers differ, and why.
          \end{itemize}
        \item \hl{Reading:} Bayes Rules!\ Chapters 1, 2, and 4.
        \item \hl{Homework:}
    \end{itemize}
\end{frame}


%%%%%%%%%%%%%%%%%%%%%%%%%%%%%%%%%%%%
% Learning objectives:
%%%%%%%%%%%%%%%%%%%%%%%%%%%%%%%%%%%%
\begin{frame}
    \frametitle{Lecture Objectives}
    \goals{%
    \begin{itemize}
    \item compute a conditional probability and apply \hl{Bayes' rule} to reverse it, using the law of total probability for the denominator;
    \item explain why a positive test for a \hl{rare} condition can still leave the condition unlikely, because the prior matters;
    \item build a probability model for a \hl{discrete parameter}: put a prior on it, evaluate the likelihood, and normalize to get the posterior;
    \item update \hl{sequentially} as data arrives, and explain why the final posterior does not depend on the order;
    \item contrast the frequentist and Bayesian answers to the same question, and say which one each actually answers.
    \end{itemize}}
\end{frame}


%%%%%%%%%%%%%%%%%%%%%%%%%%%%%%%%%%%%
\begin{frame}
    \frametitle{Lecture Objectives}
    \objectives{%
    \begin{itemize}
    \item \textbf{(M4, LO1)} Bayesian Inference
    \item \textbf{(M4, LO2)} Comparing Bayesian and Frequentist Approaches
    \end{itemize}}
    \vspace{0.3em}
    {\small Building on, and revisiting:}
    \objectives{%
    \begin{itemize}
    \item \textbf{(M2, LO1)} Approaches to Inference: p-values and s-values from Lecture 8, now set against the Bayesian answer to the same question
    \end{itemize}}
\end{frame}
